\title{Byte Latent Transformer: Patches Scale Better Than Tokens}

\begin{abstract}
We present the Byte Latent Transformer (\textsc{BLT}), a novel large language model architecture that operates at the byte level. For the first time, \textsc{BLT} achieves parity in performance with LLMs that use conventional tokenization, while delivering marked advances in inference robustness and efficiency. 

\textsc{BLT} transforms input bytes into patches of varying length, which function as the core computational units. This patching is dynamically guided by measuring the entropy of upcoming bytes, assigning increased model capacity and computation to segments exhibiting greater data complexity. We perform the inaugural controlled scaling investigation of byte-level models based on floating point operations, scaling up to 8B parameters and training on 4T bytes. Our findings underscore the practicality of training byte-based architectures at scale without relying on a predefined vocabulary. Dynamically preferring longer patches where the data exhibits regularity not only boosts training and inference throughput, but also leads to qualitative gains in complex reasoning and improved generalization across rare or intricate examples. In sum, at constant inference budgets, \textsc{BLT} achieves much better scaling behavior than its token-based counterparts by enabling simultaneous expansion of both patch lengths and model dimensionality.
\end{abstract}

\section{Introduction}
\label{section:intro}

We present the Byte Latent Transformer~(\textbf{BLT}), an architecture that operates directly on raw byte sequences, eliminating the need for tokenization. For the first time, such a tokenizer-free approach achieves parity in performance with token-based models when scaled appropriately, while delivering notable gains in both robustness and computational efficiency (\S\ref{sec:robustness}). Conventional large language models (llms) typically rely on end-to-end training, but retain an exception: tokenization, a heuristic pre-processing procedure that segments byte streams into a predefined, immutable set of tokens. This token-based compression introduces various biases, which result in limitations including sensitivity to domain and modality shifts~\citep{dagangetting}, increased vulnerability to input noise~(\S\ref{sec:robustness}), restricted knowledge of orthography~\citep{edman2024cute}, as well as inequities across languages~\citep{liang2023xlm,petrov2024language,limisiewicz2024myte}.

Historically, tokenization has played a vital role because training llms directly on byte sequences incurs excessive computational expense, largely due to the increased sequence lengths~\citep{xue2022byt5}. To address this challenge, earlier research has leveraged either more efficient variants of self-attention~\citep{el2020characterbert, clark2022canine} or architectures that eliminate attention altogether~\citep{wang2024mambabyte}~(\S\ref{section:related_work}). Nonetheless, these approaches are mainly advantageous for training \textit{small models}. For large-scale models, the primary computational bottleneck arises from the substantial feed-forward network layers applied to every byte, rather than from the self-attention mechanism.

To optimize computational resources, we introduce an adaptive, trainable approach for assembling bytes into \textit{patches}~(\S\ref{section:patching}), along with a novel model architecture that integrates both byte-level and patch-level representations. In contrast to tokenization, {\textsc BLT} does not rely on a predefined patch vocabulary. Instead, it utilizes efficient, learnable encoder and decoder modules to transform arbitrary collections of bytes into latent patch embeddings. Our results demonstrate that this strategy enables \textit{greater} computational efficiency compared to traditional tokenization-based architectures.

Tokenization-based language models assign an equal amount of computation to each token, regardless of their individual complexity. This approach prioritizes computational efficiency but can be suboptimal for performance, as the compression heuristics used in token induction do not necessarily correspond to the actual difficulty of predicting specific tokens. Our model takes a different approach by introducing mechanisms for dynamically distributing computational resources based on need. For instance, substantial model capacity is generally unnecessary for generating predictable parts of language, such as word endings, which involve lower-entropy choices compared to more challenging tasks like selecting the initial word of a sentence. This design philosophy underpins the structure of {\textsc BLT}~(\S\ref{section:architecture}), which incorporates three distinct transformer modules: two compact, byte-level \textit{local models}, and a larger, more powerful \textit{latent transformer} responsible for global reasoning~(\autoref{fig:arch_high_level}).

{\textsc BLT} determines the segmentation of bytes into patches—and thus the allocation of computational effort—by analyzing the entropy associated with predicting the next byte. This results in context-dependent patches where the information density is approximately uniform, enabling the model to concentrate compute where the predictive uncertainty, and thus the modeling challenge, is greatest.

In this work, we conduct the first extensive flop-controlled scaling investigation of byte-level models, reaching up to 8B parameters and processing 4T training bytes. Our findings demonstrate that it is feasible to train models of this scale entirely end-to-end from bytes, dispensing with the need for predetermined tokenization schemes. Notably, {\textsc BLT} achieves comparable training flop-controlled effectiveness\footnote{We assess model computational requirements by quantifying the Floating Point OPerations (flops) executed.} to Llama 3, while reducing inference flops by as much as 50\%~(\S\ref{section:scaling}). 

Moreover, directly leveraging raw byte inputs leads to substantial gains in modeling infrequent and rare patterns present in data. We find {\textsc BLT} models exhibit superior robustness to input noise compared to models utilizing tokenization, and demonstrate improved capabilities in processing at the character level, as evidenced by tasks spanning orthographic understanding, phonological analysis, and low-resource machine translation~(\S\ref{sec:robustness}).

Lastly, {\textsc BLT} enables concurrent scaling of both model and patch sizes without exceeding a predetermined inference flop limit. Using longer patch sizes tends to reduce the number of times the global latent transformer must be applied, thereby saving computational resources that can be used to increase model capacity. Our inference-flop controlled scaling studies~(\autoref{fig:fixed_inference_scaling}) reveal consistently more favorable scaling dynamics in comparison to conventional tokenization-dependent architectures.

To conclude, our work contributes the following: 1) We present {\textsc BLT}, a byte latent LLM architecture designed to enhance FLOP efficiency through dynamic compute allocation; 2) We demonstrate that, up to the 8B parameter scale, training with FLOP-matched budgets achieves comparable results to Llama 3, with the possibility of gaining up to a 50\% increase in FLOP efficiency in exchange for only marginal decreases in evaluation performance; 3) {\textsc BLT} introduces a novel approach to scaling LLMs, enabling increases in model capacity while holding the inference cost constant; 4) We provide evidence that {\textsc BLT} exhibits greater resilience to noisy inputs and effectively captures sub-word information in input sequences, an area where token-based LLMs often fall short. The codebase for training and inference with {\textsc BLT} is available at~\url{https://github.com/facebookresearch/blt}.

\section{Patching: From Individual Bytes to Groups of Bytes}
\label{section:patching}

Dividing bytes into \textit{patches} enables {\textsc BLT} to flexibly allocate computational resources depending on the surrounding context. Figure~\ref{fig:patching_types} illustrates several approaches for grouping bytes into patches. More precisely, a patching function $f_p$ partitions a byte sequence $\pmb{x} = \{x_i | i = 1, \ldots, n\}$ of length $n$ into $m < n$ patches $\pmb{p} = \{p_j | j = 1, \ldots, m\}$, assigning each $x_i$ a value in the set \{0,1\}; a value of 1 marks the start of a new patch. In both token-based and patch-based architectures, the principal determinant of computational cost is the number of operations required by the core Transformer, which, for {\textsc BLT}, is dictated by how many patches are generated from the data via the chosen patching function. As such, the average number of bytes per patch—or simply, the \textit{patch size}—emerges as the key variable influencing processing cost during both training and inference under a specified patching strategy~(\S\ref{section:flops}). 

We proceed by introducing three specific patching methodologies: segmentation into patches containing a fixed byte count~(\S\ref{section:static-patch}), patching determined by whitespace boundaries~(\S\ref{section:white-patch}), and a dynamic scheme that leverages entropy measures from a compact byte language model~(\S\ref{section:dyn-patch}). We conclude by exploring incremental patching, and examining distinctions between tokenization and patching~(\S\ref{section:bpe-patch}).

\subsection{Strided Patching Every K Bytes}
\label{section:static-patch}

A common approach for grouping bytes involves segmenting them into fixed-size patches of length $k$, as utilized in MegaByte~\citep{yu2023megabyte}. This method, based on a constant stride, offers ease of implementation during both training and inference. It also allows for simple adjustment of the mean patch size, enabling straightforward manipulation of the computational cost (flop cost). Nevertheless, this strategy is accompanied by notable disadvantages. Primarily, it lacks adaptive computation—resources are not directed to the most information-rich regions: a transformer layer $j$ might be underutilized when predicting uninformative regions such as whitespace, or conversely, regions dense in complex content (e.g., mathematical notation) might not receive adequate attention. Additionally, this approach can cause similar sequences of bytes, like a repeated word, to be segmented inconsistently and without consideration of their contextual meaning.

\subsection{Space Patching}
\label{section:white-patch}

\citet{slagle2024spacebyte} introduces a straightforward yet powerful modification to strided patching, in which new patches are created after encountering space-like bytes\footnote{
    Space-like bytes are defined as any byte that is not a latin character, digit, or utf-8 continuation byte. In addition, each patch must contain at least one non space-like byte. }, which frequently serve as natural delineations for linguistic units in various languages. In the Space patching approach, each word receives a dedicated latent transformer computation step (entailing additional flops), ensuring consistent patching of words across different sequences and focusing computational resources on challenging predictions, which typically arise after spaces. For instance, when completing the question ``Who composed the Magic Flute?\uline{\hspace{.6em}}'', predicting the initial byte of the answer is far more difficult than generating subsequent bytes following ``M'', since the initial character drastically narrows the set of plausible continuations, thereby making the completion ``Mozart'' relatively straightforward to predict. Nonetheless, space patching is not universally applicable across all languages and domains and, crucially, lacks flexibility in controlling patch size. In the subsequent section, we present a new patching technique inspired by the observation that the leading bytes of words usually pose the greatest challenge for prediction, while offering a natural means of varying patch size.

\subsection{Entropy Patching: Using Next-Byte Entropies from a Small Byte LM}
\label{section:dyn-patch}

Instead of utilizing a heuristic rule like whitespace for segmentation, we adopt a data-centric methodology to pinpoint next-byte predictions with elevated uncertainty. We propose \textit{entropy patching}, a technique that employs entropy estimation to determine the locations of patch boundaries.

We train a small byte-level auto-regressive language model on the training data for {\textsc BLT} and compute next byte entropies under the LM distribution $p_{e}$ over the byte vocabulary $\mathcal{V}$:

\begin{equation}
H(x_i) = \sum_{v \in \mathcal{V}} p_{e}(x_i=v|\pmb{x}_{<i}) \log p_{e}(x_i=v|\pmb{x}_{<i})
\end{equation}

We explore two approaches for determining patch boundaries based on entropies $H(x_i)$. The first method involves selecting points whose entropy exceeds a fixed global threshold, as shown in~\autoref{fig:patching}. The second strategy targets points where the entropy is notably higher compared to the immediately preceding value. This latter method can also be viewed as detecting points that disrupt an otherwise roughly monotonically decreasing entropy sequence within a patch.

\begin{align*}
    \text{Global Constraint}\hspace{1cm}H(x_t)& > \theta_g\\
    \text{Approx. Monotonic Constraint}\hspace{1cm}H(x_t)& - H(x_{t-1}) > \theta_r
\end{align*}

Patch boundaries are determined through an efficient preprocessing phase performed as part of dataloading. This approach contrasts with~\citet{nawrot-etal-2023-efficient}, who employ a classifier trained to detect entropy-based patch boundaries. In our experiments~(\S\ref{section:experiments}), we evaluate and compare these two strategies for identifying bytes with low versus high entropy.

\subsection{The Byte-Pair Encoding (BPE) Tokenizer and Incremental Patching}
\label{section:incremental-patching}
\label{section:bpe-patch}

A wide range of contemporary llms, such as our baseline Llama 3, employ subword tokenizers, for instance bpe~\citep{gage1994new, sennrich-etal-2016-neural}. Throughout this work, we define ``tokens'' as groups of bytes selected from a \textit{finite} vocabulary that is established before training, contrasting with ``patches'', which are sequences formed dynamically and lack a pre-defined vocabulary. A key distinction is that, in the case of tokens, the model is not provided with explicit access to the original byte-level features.

A significant advancement that {\textsc BLT} introduces compared to models based on tokenization lies in its redefinition of the balance between vocabulary size and computational expense. In conventional large language models (LLMs), expanding the vocabulary leads to tokens that typically encompass more bytes, which results in fewer processing steps per input; however, this also increases the dimensionality of the model’s output layer, particularly in the final projection. This inherent compromise greatly limits tokenization-centric models from substantially modifying both average token length and inference requirements. As an illustration, Llama 3 raises its average token length from 3.7 to 4.4 bytes, but this enhancement comes at the price of expanding its embedding matrix by a factor of four relative to Llama 2.

During generation, {\textsc BLT} must determine at each position in the byte sequence whether a patch boundary has been reached, since this dictates whether additional computation through the Latent Transformer is required. This boundary detection must be accomplished without reliance on any bytes that have yet to be generated; in other words, the segmentation process must not depend on future context. Formally, an incremental patching scheme $f_p$ must fulfill the following criterion:
$$f_p(\pmb{x}_{<i}) = f_p(\pmb{x})_{<i}$$
bpe does not qualify as an incremental patching scheme, because a given prefix may receive different tokenizations based on what follows it, violating the above condition\footnote{Introducing a special delimiter token at patch boundaries could make bpe comply with incremental patching, but it would lengthen the resulting byte sequence.}.

\section{{\textsc BLT} Architecture}
\label{section:architecture}
The {\textsc BLT} framework consists of a primary global autoregressive language model working on patch-level representations, supplemented by two more compact local modules. These local models are responsible for transforming byte sequences into patches for input and reconstructing bytes from patch representations as output~(\autoref{fig:arch_high_level}).

\subsection{Latent Global Transformer Model}

\textit{The Latent Global Transformer} serves as an autoregressive transformer model $\mathcal{G}$ composed of $l_{\mathcal{G}}$ layers. This model translates a sequence of latent input patch embeddings, $p_j$, into corresponding output patch embeddings, $o_j$. In this work, patches are indexed by the subscript $j$, while the subscript $i$ refers to bytes. The global model leverages a block-causal attention mask~\citep{dubey2024llama}, enforcing the attention mechanism to consider only patches up to and including the current patch within a document. Notably, this model is responsible for the majority of computational cost both during pre-training and inference. As such, selectively invoking this model enables adaptive allocation of computational resources based on the complexity of specific input and output segments.

\subsection{Local Encoder}
\label{section:local}

The \textit{Local Encoder Model}, represented as $\mathcal{E}$, comprises a compact transformer architecture with significantly fewer layers ($l_{\mathcal{E}} << l_{\mathcal{G}}$) compared to the global model. Its principal function is to rapidly encode an input sequence of bytes $b_i$ into a set of informative patch-level representations $p_j$. A key distinction from standard transformer models lies in the integration of a cross-attention mechanism after each transformer layer, which serves to aggregate byte-level embeddings into corresponding patch-level representations~(\autoref{fig:crossattn}).

Initially, each input byte $b_i$ is projected into an embedding space via a matrix of size $\mathbb{R}^{256 \times h_{\mathcal{E}}}$, producing embeddings $x_i$. Optionally, these embeddings can be further enhanced using hash-embeddings, as elaborated in~(\S\ref{sec:hash-emb}). Through a sequence of transformer and cross-attention modules~(\S\ref{sec:enc-cross-attn}), the model progressively refines these embeddings, culminating in the patch representations $p_i$, which are subsequently handed off to the global transformer $\mathcal{G}$ for further processing.

The attention mechanism within each transformer block employs a \textit{local block causal} masking strategy: each byte may attend to up to $w_{\mathcal{E}}$ bytes that precede it, which may span across flexible patch delineations but will not breach the confines of the original document. The subsequent sections provide more in-depth discussion of the embedding methods and the inner workings of the cross-attention component.

\subsubsection{Encoder Hash n-gram Embeddings}
\label{sec:ngram-emb}
\label{sec:hash-emb}
An essential aspect of forming strong and informative representations at every step $i$ involves embedding contextual details regarding the bytes that come before. Within {\textsc BLT}, this is accomplished by representing the current byte $b_i$ both on its own and as a member of a larger byte n-gram. At every step $i$, the model first builds byte-grams

\begin{align}
    g_{i,n}=\{b_{i-n + 1},\ldots, b_i\}
\end{align}

for every byte position $i$ and for values of $n$ ranging from three up to eight.\footnote{
Byte-grams of size $n$ or larger are omitted in cases where $i<n$. }

Next, we propose hash $n$-gram embeddings, which utilize a hash function to project all byte $n$-grams to an index within a fixed-size embedding matrix $E_{n}^{hash}$ for each $n \in\{3,4,5,6,7,8\}$~\citep{bai2010learning}. This embedding is subsequently combined with the original byte embedding, followed by normalization, before serving as input to the local encoder module. The enriched embedding is computed as follows:

\begin{align}
e_i &= x_i + \sum_{n = 3,...,8}E_{n}^{hash}(\text{Hash}(g_{i,n})) \\
\text{where, Hash}(g_{i,n}) &= \text{RollPolyHash}(g_{i,n}) \% |E_{n}^{hash}|
\end{align}

We scale $e_i$ by dividing by the total number of $n$-gram lengths considered plus one, and implement RollPolyHash as specified in Appendix~\ref{appendix:rollpolyhash}. Section~\ref{section:ablations} presents an ablation study examining the impact of varying $n$-gram hash embedding parameters, such as $n$ values and embedding table sizes, on the observed scaling law patterns under fixed computational budgets. Beyond utilizing hashed $n$-gram embeddings, our investigations extended to frequency-based $n$-gram embeddings; a comprehensive account of these experiments is available in Appendix~\ref{appendix:ngrams}.

\subsubsection{Encoder Multi-Headed Cross-Attention}
\label{sec:enc-cross-attn}
Our approach is largely inspired by the input cross-attention mechanism from the Perceiver architecture~\citep{jaegle2021perceiver}. The primary distinction in our setup is that the latent representations here are aligned with patch representations that vary per input, rather than relying on a predefined set of latent representations~(\autoref{fig:crossattn}). Additionally, each latent only interacts with the bytes specific to its corresponding patch. Within the module, each patch $p_j$ is represented by a query vector. This vector is initialized by applying a pooling operation over the byte embeddings associated with patch $p_j$, and is subsequently passed through a linear transformation, $\mathcal{E}_{C} \in \mathbb{R}^{h_{\mathcal{E}} \times (h_{\mathcal{E}} \times U_{\mathcal{E}})}$, in which $U_{\mathcal{E}}$ denotes the quantity of encoder cross-attention heads. Explicitly, letting $f_{\text{bytes}}(p_j)$ signify the sequence of bytes comprising patch $p_j$, we compute

\begin{align}
P_{0,j} &= \mathcal{E}_{C}(f_\text{bytes}((p_j)), f~ \text{is a pooling function} \\
P_l &= P_{l-1} + W_o\left(\text{softmax}\left(\frac{QK^T}{\sqrt{d_k}}\right)V\right) \\
\text{where } Q_j &= W_q(P_{l-1,j}), K_i = W_k(h_{l-1,i}), V_i = W_v(h_{l-1,i}) \\
h_l &= \text{Encoder-Transformer-Layer}_l(h_{l-1})
\end{align}

Here, $P \in \mathbb{R}^{n_p \times h_{\mathcal{G}}}$ denotes the collection of $n_p$ patch representations meant for input to the global model, initialized by aggregating the byte embeddings $e_i$ associated with each patch $p_j$. The matrices $W_q$, $W_k$, $W_v$, and $W_o$ are designated as the linear projection matrices for queries, keys, values, and output, respectively, where keys and values are derived from the byte-level representations $h_i$ output by the preceding layer (or $e_i$ for the first layer). A tailored masking method is applied in which each query $Q_j$ exclusively attends to keys and values that are associated with bytes within the same patch $j$. Multi-head attention is performed over $Q$, $K$, and $V$. Given that patch representations have a higher dimensionality ($h_{\mathcal{G}}$) than $h_{\mathcal{E}}$, $P_l$ is organized as multiple heads of size $h_{\mathcal{E}}$ during cross-attention and subsequently concatenated to recover the $h_{\mathcal{G}}$-dimensional representation. Furthermore, pre-LayerNorm is applied separately to the queries, keys, and values, and no positional embeddings are incorporated into this cross-attention component. A residual connection is applied to encompass the cross-attention module.

\subsection{Local Decoder} 

The local decoder $\mathcal{D}$, mirroring the structure of the local encoder, is designed as a compact transformer model comprised of $l_{\mathcal{D}} << l_{\mathcal{G}}$ layers. Its primary function is to convert a sequence of global patch representations $o_j$ back into the original sequence of raw bytes $y_i$. Tasked with predicting each raw byte in sequence, conditioned on all previously generated bytes, it utilizes the hidden representations produced by the local encoder for the byte sequence as input. The local decoder processes data through $l_{\mathcal{D}}$ layers that alternate between cross-attention and standard transformer operations. Specifically, the cross-attention mechanism is applied prior to each transformer layer, enabling the creation of byte representations from patch-level features, while the subsequent transformer layers refine these representations along the decoded byte sequence.

\subsubsection{Decoder Multi-headed Cross-Attention} 

In the decoder cross-attention, the roles of the queries and key/values are interchanged i.e. the byte-representations are now the queries, and the patch representations are now the key/values. The initial byte-representations for the cross-attention are initialized as the byte embeddings from the last encoder layer i.e. $h_{l_{\mathcal{E}}}$. The subsequent byte-representations for layer $l$, $d_{l,i}$ are computed as:

\begin{align}
D_0 &= h_{l_{\mathcal{E}}} \\
B_l &= D_{l-1} + W_o\left(\text{softmax}\left(\frac{QK^T}{\sqrt{d_k}}\right)V\right), \\
  &\text{where } Q_i = W_q(d_{l-1,i}), K_i = W_k(\mathcal{D}_C(o_j)), V_i = W_v(\mathcal{D}_C(o_j)) \\
  D_l &= \text{Decoder-Transformer-layer}_l(B_l)
\end{align}

Here, as previously noted, $W_k$ and $W_v$ denote key and value projection matrices, respectively; these apply a linear mapping along with the split operation $\mathcal{D}_C$ to the final patch-level outputs $o_j$ originating from the global model. The matrix $W_q$ is utilized to project queries, operating on the byte-level activations $d_{l-1}$ passed from the preceding decoder transformer layer (for the initial decoder layer, $h_{l_{\mathcal{E}}}$ is used instead). $W_o$ serves as the output projection matrix. Consequently, $B$ has dimensions $\mathbb{R}^{h_{\mathcal{D}} \times n_b}$, where $n_b$ corresponds to the total number of bytes in the output. The subsequent decoder layer representations, $D_l$, are then generated via a decoder transformer layer fed with the cross-attention output $B$. Analogous to the setup in the local encoder’s cross-attention, our architecture incorporates multi-head attention, employs pre-LayerNorm, excludes positional embeddings, and encloses the cross-attention block within a residual connection.

\section{Experimental Setup}
\label{section:experiments}
We have meticulously established a controlled experimental framework to conduct a rigorous comparison between {\textsc BLT} and models that employ tokenization, making certain that {\textsc BLT} does not benefit from potentially increased sequence context lengths.

\subsection{Pre-training Datasets} In our experiments, all model scales are pre-trained using two primary datasets: 1) the Llama 2 dataset~\citep{touvron2023llama}, which consists of 2 trillion tokens sourced from diverse publicly accessible materials, later subjected to extensive cleaning and filtering for quality enhancement; and 2) {\textsc BLT}-1T, a new dataset containing 1 trillion tokens gathered from a range of public sources, which also incorporates a subset of the Datacomp-LM pre-training corpus~\citep{li2024datacomp}. The Llama 2 dataset supports scaling law investigations on the optimal token count, following the methodology of~\cite{dubey2024llama}, aiding in the selection of effective architectural strategies for {\textsc BLT}. In contrast, {\textsc BLT}-1T enables a comprehensive pre-training cycle intended for downstream task comparisons with Llama 3. It is important to note that neither dataset includes material obtained from Meta products or services. Additionally, for the tokenizer-based baseline studies, the Llama 3 tokenizer—featuring a 128K vocabulary—was adopted, as it demonstrated stronger results in our benchmarks relative to the Llama 2 tokenizer.

\subsection{Entropy Model} For our experiments, the entropy model is implemented as a byte-level language model, which is trained using the same data distribution as the complete {\textsc BLT} model. By default, unless otherwise specified, the architecture consists of a transformer comprising 100 million parameters, 14 layers, a hidden size of 512, and employs sliding window attention across 512 bytes. All other hyperparameters mirror those utilized in both our local and global transformer setups. We have also evaluated alternative model capacities, receptive fields, and architectural variants as detailed in \autoref{section:ablations}. Notably, when the model's receptive field is reduced sufficiently, the resultant entropy model can be effectively represented using a compact lookup table.

\subsection{Entropy Threshold and Equalizing Context Length} For models that implement entropy-driven patching, we determine a threshold for patching intended to produce a specified mean \textit{patch size} across the pretraining dataset mixture. In {\textsc BLT}, in contrast to standard tokenization, the selection of \textit{patch size} is fully flexible, which substantially influences the model's accessible context window. To prevent models with larger patch sizes from gaining an undue benefit due to longer contexts, we ensure that the expected total byte count per batch is held constant. As a result, when employing larger patch sizes, we correspondingly decrease the number of sequence steps. Specifically, for Llama 2 training data, an 8k byte context window is used; for the {\textsc BLT}-1T dataset, the context is raised to 16k bytes on average, while the average batch size remains fixed at 16M bytes.

Although the average batch size remains unchanged, dynamic patching methods produce varying byte-to-patch ratios across data batches. For the sake of efficiency, our implementation of {\textsc BLT} training organizes patch batches in a manner that eliminates the need for padding operations within the computationally intensive latent transformer, thereby guaranteeing a consistent patch count per batch. To prevent memory usage surges caused by atypically large patches, we pad and, where necessary, truncate input byte sequences to fixed lengths of 12k and 24k bytes for the Llama 2 and {\textsc BLT}-1T datasets, respectively, during training.

\subsection{Entropy Model Context}
\label{section:entropy-context}
Our empirical observations indicate that entropy patching leads to increasingly larger patches, particularly in highly structured content such as multiple choice formats, which tend to be quite repetitive (see the example in \autoref{fig:mmlu_patching}). This expansion is primarily a consequence of the lower entropy associated with the redundant elements present within the entropy model context. Therefore, during the large-scale application of {\textsc BLT}-Entropy with a patch size of 4.5, we opt to reset the entropy context with line breaks and employ an approximate monotonicity constraint, as it shows greater resistance to "entropy drift" that arises due to variations in context length. While this adjustment influences only our entropy calculation method, our approach to determining the appropriate entropy threshold remains unchanged.

\subsection{FLOPs Estimation} 
\label{section:flops}

Our approach closely mirrors the methodology for calculating transformer flops as presented in Chinchilla~\citep{hoffmann2022training}, accounting for the operations associated with the feed-forward components, qkvo projections within the self-attention mechanism, as well as the steps for attention computation and output projection. However, an important distinction in our work is the treatment of the input embedding layer: we assume it utilizes an efficient lookup rather than performing dense matrix multiplication, and thus we attribute zero flops to this stage. Consistent with prior literature, we approximate the computational cost of the backward pass as double that of the forward pass.

To compute flops \textit{per byte} for {\textsc BLT} models, we add up the flops for the local encoder transformer, the global latent transformer, and the local decoder transformer, together with the cross attention blocks in the encoder and the decoder:

\begin{align}
    \text{FL}_{\text{{\textsc BLT}}} &= \text{Transf. FL}(h_{\mathcal{G}}, l_{\mathcal{G}}, m=n_{ctx}/n_p, V=0) / n_p\\
   &+ \text{Transf. FL}(h_{\mathcal{E}}, l_{\mathcal{E}}, m=w_{\mathcal{E}}, V=0) \\
   &+ \text{Transf. FL}(h_{\mathcal{D}}, l_{\mathcal{D}}, m=w_{\mathcal{D}}, V=256) \\ 
    &+ \text{Cross Attn. FL}(h_{\mathcal{E}}, l_{\mathcal{E}}, m=n_p, r=n_p/k) \cdot \frac{k}{n_p} \\ 
   &+ \text{Cross Attn. FL}(h_{\mathcal{D}}, l_{\mathcal{D}}, m=k, r=k/n_p) 
\end{align}

Here, $n_{ctx}$ denotes the sequence length measured in bytes, while $n_p$ specifies the patch size. The parameter $r$ captures the proportion of queries relative to keys and values, and $k$ defines the ratio between the patch dimension and the byte dimension—that is, it represents the count of local model segments merged to construct a comprehensive global representation (for instance, $k=2$ as shown in \autoref{fig:crossattn}). The term $V$ refers to the vocabulary size used solely within the local decoder’s output projection. The attention module employs varying context lengths, $m$, depending on whether it operates over bytes or patches. Consequently, we adjust the flop computations for attention to reflect these differences for each respective module. The precise formulations for calculating Transformer-FLOPs and Cross-Attention FLOPs are detailed in Appendix~\ref{section:flops-appendix}.

\subsection{Bits-Per-Byte Estimation} Perplexity only makes sense in the context of a fixed tokenizer as it is a measure of the uncertainty for each token. When comparing byte and token-level models, following previous work~\citep{xue2022byt5, yu2023megabyte, wang2024mambabyte}, we instead report Bits-Per-Byte (BPB), a tokenizer independent version of perplexity. Specifically:

In this expression, the level of uncertainty associated with the dataset $\pmb{x}$, quantified using the aggregate cross-entropy loss, is divided by both the total byte count of $\pmb{x}$ and a scaling constant in order to normalize the measure.

\subsection{Transformer Architecture Hyperparameters}  
Across all transformer modules in {\textsc BLT}, encompassing both local and global components, we adopt design choices similar to those implemented in Llama~3~\citep{dubey2024llama}. Specifically, our feed-forward networks utilize the SwiGLU activation~\citep{swiglu}, and rotary positional embeddings (RoPE)~\citep{RoPe} are employed in the self-attention layers with a parameter setting of $\theta=500000$~\citep{xiong2024effective}. For normalization, we rely on RMSNorm~\citep{RMSNorm}. Fixed-standard attention masks—such as \textit{block causal} or \textit{fixed-window block causal}—in self-attention layers are handled using Flash Attention~\citep{dao2022flashattention}, and for fixed-width masks, we set the window size to 512. When dealing with dynamic, patch-specific masks within cross-attention modules, we adopt Flex Attention\footnote{\url{https://pytorch.org/blog/flexattention}}, which delivers efficient fused operations that considerably reduce training times.

\subsection{{\textsc BLT}-Specific Hyperparameters} In order to evaluate the performance of {\textsc BLT} models, our experimental design encompasses two main aspects: examining scaling behaviors and assessing downstream tasks. For these analyses, we train models of various sizes: 400M, 1B, 2B, 4B, and 8B parameters. Detailed architectural hyperparameters for these models are provided in Appendix~Table~\ref{tab:arch}.

Initialization of queries for the initial cross-attention layer in the local encoder is performed via max-pooling. All variants of {\textsc BLT} utilize $500,000$ hashes and a single hash function, considering n-gram lengths from 3 to 8. The learning rate is set uniformly to $4e-4$ across all configurations. This learning rate selection aligns with results from a sweep over $1e-3$ to $1e-4$ at the 400M and 1B scales, which indicated that the same rate is optimal for both token-based and {\textsc BLT} models. For scaling experiments using Llama-2 data, the batch sizes adhere to those recommended in~\cite{dubey2024llama}, or are equivalently matched by data volume in bytes. Model optimization employs the AdamW algorithm~\citep{adamw}, setting $\beta_1$ to 0.9, $\beta_2$ to 0.95, and $\epsilon=10^{-8}$. We employ a linear warm-up schedule for 2000 steps, followed by cosine annealing to decay the learning rate to zero. Furthermore, a weight decay factor of 0.1 is applied, alongside global gradient clipping set at 1.0.

\section{Scaling Trends}
\label{section:scaling}

We offer a comprehensive overview of the scaling behavior exhibited by byte-level models, providing insights relevant to the continued scaling of {\textsc BLT} architectures. This scaling analysis is designed to overcome shortcomings of prior studies involving byte-level models by: (a) systematically analyzing trends within a compute-optimal training paradigm; (b) training 8B parameter models using substantial datasets (up to 1T tokens or 4T bytes) and assessing their effectiveness on various downstream applications; and (c) quantifying scaling effects under fixed inference-cost conditions. In subsequent sections, we will delve into the distinct benefits arising from direct modeling of byte sequences.

\subsection{Parameter Matched Compute Optimal Scaling Trends}
\label{section:parameter-matched-scaling}
Leveraging the Llama 2 dataset, we conduct training runs for both \textit{compute-optimal} bpe models and {\textsc BLT} models at four parameter scales, spanning from 1B to 8B. We analyze their training efficiency by charting the total training flops against language modeling outcomes, utilizing a representative segment of the training data mixture. For bpe models, training follows the parameter-to-data ratio identified as optimal by Llama~3~\citep{dubey2024llama}. Such a \textit{compute-optimal} configuration, as supported by theoretical analyses~\citep{hoffmann2022training}, aims to maximize performance for a fixed computational budget, thereby forming a reliable baseline. To ensure fair comparison, we pair each bpe model with a corresponding {\textsc BLT} model, applying identical data and employing a Latent Transformer architecture with matching scale and design.

As shown in~\autoref{fig:scaling-compute-opt} (right), {\textsc BLT} models consistently achieve results that are on par with, or surpass, those of their bpe counterparts, a pattern that persists as both model size and computational effort increase. To our knowledge, this work represents the first instance where a byte-level Transformer architecture like {\textsc BLT} demonstrates scaling behaviors that closely align with BPE-based models under compute-optimal conditions. This observation provides evidence supporting our hypothesis that the parameter-to-compute ratio considered optimal for bpe models is also largely applicable to {\textsc BLT} architectures.

Achieving parity with bpe scaling patterns requires advancements both in architectural design and dynamic patching. In ~\autoref{fig:scaling-compute-opt} (left), we present a comparison between models leveraging space-patching techniques and Llama 3. To represent SpaceByte \citep{slagle2024spacebyte}, we use {\textsc BLT} configured for space-patching, while omitting n-gram embeddings and cross-attention mechanisms. While SpaceByte surpasses Megabyte in performance, a significant gap persists between its results and those of Llama 3. In \autoref{fig:scaling-compute-opt} (right), we display the cumulative gains derived from both improved architectures and dynamic patching methods. When scaling up, {\textsc BLT} models achieve competitive results, closely matching state-of-the-art tokenizer-based systems like Llama 3.

Additionally, we note that the selection of tokenizer significantly influences performance in tokenizer-based models; specifically, models utilizing the Llama-3 tokenizer demonstrate superior performance compared to those employing the Llama-2 tokenizer when both are trained on identical datasets.

In summary, when utilizing much larger patch sizes, the {\textsc BLT} architecture demonstrates scaling behavior that sits between Llama 2 and Llama 3. The Llama 2 and 3 models, which use bpe tokenizers, display mean token lengths of 3.7 and 4.4 bytes, respectively. By contrast, {\textsc BLT} maintains comparable scaling characteristics when the mean patch size is raised to 6 or even 8 bytes. Since inference flops scale inversely with average patch size, increasing the patch size to 8 bytes results in nearly a 50\% reduction in required inference flops. Notably, models configured with larger patch sizes exhibit improved performance as both model scale and data volume are increased. Although {\textsc BLT} using an 8-byte patch size starts from a lower performance baseline at the 1B parameter scale compared to bpe-tokenized Llama 2, it surpasses the latter at 7B parameters. This pattern implies that further increases in patch size could yield additional benefits at even greater model sizes and computational budgets, suggesting the viability of experimenting with even larger patch sizes as resources allow.

\subsection{Beyond Compute Optimal Task Evaluations}
\label{section:evals}
To further investigate scaling dynamics, we extend training of an 8B {\textsc BLT} model past the compute-optimal regime using the {\textsc BLT}-1T dataset, which consists of larger and higher-quality data, and systematically measure the model’s capabilities on various established classification and generation benchmarks. We focus our evaluation on tasks pertaining to common sense reasoning, world knowledge, and code generation:
\paragraph{Classification tasks} encompass ARC-Easy (0-shot)~\citep{Eval_ARC}, Arc-Challenge (0-shot)~\citep{Eval_ARC}, HellaSwag (0-shot)~\citep{Eval_hellaswag}, PIQA (0-shot)~\citep{Eval_piqa}, and MMLU (5-shot)~\citep{hendrycks2020measuring}. For these, we adopt a prompt-based scoring approach, computing the probabilities assigned to different answer choices, and present average accuracy scores.
\paragraph{Coding related generation tasks:} We evaluate the models’ proficiency in Python code generation by reporting pass@1 performance on MBPP (3-shot)~\citep{Eval_MBPP} and HumanEval (0-shot)~\citep{Eval_HumanEval}.

In \autoref{tab:evals}, we present a comparison among three models, all trained on the {\textsc BLT}-1T dataset: a model using the Llama 3 tokenizer with byte pair encoding,\footnote{The Llama 3 tokenizer was selected for its 128k vocabulary, as it demonstrates superior performance compared to Llama 2's 32k vocabulary.} and two distinct variants of the {\textsc BLT} architecture. The first variant applies a space-patching approach ({\textsc BLT}-Space), while the second leverages an entropy-driven patching mechanism ({\textsc BLT}-Entropy) that incorporates an approximate monotonicity constraint and resets the context after new line tokens, following the protocol in~\autoref{section:entropy-context}. All models were allocated an equivalent number of floating-point operations for training. Notably, for {\textsc BLT}-Entropy, we decrease the entropy threshold from 0.6 to 0.1 during inference, a modification which enhances performance on downstream tasks, although this requires a greater number of inference steps.

The {\textsc BLT}-Entropy model demonstrates superior performance compared to the Llama 3 model on 4 out of 7 tasks, despite both being trained on an equal number of bytes. This enhancement likely stems from two main factors: (1) more effective utilization of training computation enabled by dynamic patching, and (2) the explicit modeling of information at the byte level rather than at the token level.

Conversely, {\textsc BLT}-Space lags behind the Llama 3 tokenizer in all tasks except one; nevertheless, it delivers notable savings in inference flops due to its increased average patch size of 6 bytes. By contrast, both the bpe and entropy-patching methods produce models with similar average patch sizes, approximately 4.5 bytes, when evaluated on the mixed training dataset. Given an equal training budget, the model utilizing a larger patch size is able to process 30\% more data relative to the other two approaches. This increased data throughput may, however, cause {\textsc BLT} to deviate further from the point of compute-optimality.

\subsection{Patches Scale Better Than Tokens} {\textsc BLT} models allow for joint scaling of both model parameters and patch dimensions, all while preserving a fixed computational budget for both training and inference, as well as a constant dataset size. Unlike traditional token-based approaches with a fixed vocabulary, patch-based architectures permit flexible adjustment of patch size, effectively bypassing the efficiency limitations described in Section~\ref{section:bpe-patch}. Expanding patch size reduces computational overhead, enabling those resources to be redirected towards enlarging the global latent transformer, which benefits from being invoked at a lower frequency.

We perform a fixed inference scaling analysis to evaluate whether larger models executing fewer steps on larger input patches can surpass the performance of smaller models executing more steps. Beginning with models containing 400 million and 3.6 billion parameters utilizing the Llama 2 tokenizer, we identify flop-matched counterparts using the Llama 3 tokenizer as well as {\textsc BLT}-Entropy architectures configured with mean patch sizes of 6 and 8 bytes based on the training data mix (refer to~\autoref{tab:fixed-inf-params} for specific model configurations). Notably, for models using an 8-byte patch size, we employ 3 encoder layers instead of a single one. All models are trained across different training flop allocations.

As depicted in \autoref{fig:fixed_inference_scaling}, {\textsc BLT} models exhibit superior scaling performance compared to tokenization-based counterparts across both inference flop categories. Although BPE architectures initially outperform at lower training budgets, {\textsc BLT} models soon overtake them after slightly exceeding the compute-optimal threshold. This observation suggests that allocating additional resources to pretraining can be advantageous, yielding higher performing models within the same inference constraints. The 8B model class, including Llama 3.1—which was trained on two orders of magnitude more data than what is compute-optimal for its parameter count—exemplifies this approach.

The point at which {\textsc BLT} surpasses token-based models in performance moves slightly nearer to the compute-optimal threshold with an increase in flop class, shifting from three times to about 2.5 times the compute-optimal budget. In addition, the model utilizing a patch size of 8 within the higher flop class exhibits a sharper scaling trajectory, outpacing other configurations at an earlier stage. As outlined in Section~\ref{section:parameter-matched-scaling}, scaling to larger patch sizes narrows the gap in performance with BPE models at greater model capacities. We partly ascribe this behavior to the fact that the portion of total flops attributed to the byte-level Encoder and Decoder components diminishes as these modules scale less aggressively compared to the Latent Transformer. When the overall parameter count increases twentyfold, from 400M to 8B, the number of local model parameters in {\textsc BLT} only approximately doubles. This distinction matters because augmenting patch size impacts the compute spent on the patch Latent Transformer, without influencing that of the byte-level modules. Consequently, this explains why the parameter count for {\textsc BLT}-Entropy ps=8 rose from 1.6 times to 1.7 times that of the Llama 2 model size when scaling up to the larger models.

In conclusion, our investigation into patch-length scaling reveals that the {\textsc BLT} patch-based framework benefits from concurrent increases in both patch length and model capacity, resulting in superior scaling performance. Notably, these scaling advantages not only persist but can be further enhanced as models grow larger.

\section{Byte Modeling Improves Robustness} Additionally, we assess the robustness of {\textsc BLT} relative to models based on tokens that do not incorporate explicit byte-level details, and introduce a method for augmenting pretrained token-based models with byte-level representations.
\label{sec:robustness}

\subsection{Character-Level Tasks}

Initial research into byte-level models was driven by the expectation that they would exhibit greater resilience to noise at the byte level in input data, and that they would better capture the underlying structure of tokens—a notable limitation in present tokenizer-based architectures. To assess these characteristics, we conduct supplementary experiments using benchmarks specifically designed to test resilience to noisy inputs as well as character-level sensitivity, encompassing English, multiple other languages, digits, and phonemes. The outcomes of these evaluations are reported in Table~\ref{tab:char_tasks}.

\paragraph{Noisy Data} To evaluate the robustness of {\textsc BLT} in comparison to models dependent on tokenizers, we construct noisy variants of the benchmark classification datasets outlined in Section~\ref{section:evals}. Specifically, we introduce text perturbations using five unique character-level noise methods: (a) \textit{AntSpeak}: This technique transforms the input text into uppercase and inserts spaces between every character. (b) \textit{Drop}: Randomly eliminates 10\% of the text's characters. (c) \textit{RandomCase}: Randomly assigns uppercase to half of the characters and lowercase to the other half throughout the text. (d) \textit{Repeat}: Selects 20\% of the characters at random and repeats each up to four consecutive times. (e) \textit{UpperCase}: Converts the entire text into uppercase. For each of these perturbation methods, we separately apply noise to the prompt, the completion, or both, treat each setting as a distinct task, and calculate the average performance metrics. Table \ref{tab:char_tasks} presents results for the noised version of HellaSwag~\citep{Eval_hellaswag}. The outcomes demonstrate that {\textsc BLT} consistently shows superior resilience to noise compared to tokenizer-based baselines, outperforming a baseline model trained with the same data by an average margin of 8 points, and surpassing even the Llama 3.1 model, which has been trained with a significantly larger corpus.

\paragraph{Phonology - Grapheme-to-Phoneme (G2P)} We evaluate the proficiency of {\textsc BLT} in converting sequences of graphemes (alphabetic representations of words) into their corresponding phonemic transcriptions. Table \ref{tab:char_tasks} summarizes the 5-shot performance of the G2P task utilizing Phonology Bench~\citep{suvarna2024phonologybench}. The results indicate that {\textsc BLT} achieves superior accuracy compared to the Llama 3 1T tokenizer-based baseline on this benchmark.

\paragraph{CUTE}
To measure character-level comprehension, we test {\textsc BLT} using the CUTE benchmark~\citep{edman2024cute}, which includes multiple tasks divided broadly into three areas: compositional understanding, orthographic similarity, and sequence manipulation. This evaluation is notably demanding for most models that rely on tokenization, as they often recognize how their tokens are spelled but find it challenging to leverage this knowledge for manipulating text effectively. As shown in Table~\ref{tab:char_tasks}, {\textsc BLT}-Entropy surpasses both BPE Llama 3 models by over 25 points on this benchmark. Our model is particularly strong at character manipulation, obtaining 99.9\% accuracy on both spelling-related tasks. These significant gains are observed even though {\textsc BLT} is trained on sixteen times less data than Llama 3.1, suggesting that acquiring character-level abilities is especially difficult for BPE models. Figure \ref{fig:cute_examples} highlights examples where the Llama 3 tokenizer model encounters difficulties while our {\textsc BLT} model succeeds. The only areas where BPE excels are in word deletion and insertion; such word-level modifications are less straightforward for models operating at the byte level, though the performance gap is modest. It may, in fact, be more feasible to build word-level capabilities from characters upward than the reverse. Across all tasks, we maintain a consistent evaluation protocol and use the original prompts provided by Huggingface. It is possible that prompt engineering could further improve the performance of BPE models.

\paragraph{Low Resource Machine Translation} We assess the performance of {\textsc BLT} on translation tasks involving both directions—into and out of—six widely represented language families and twenty-one low-resource languages, each employing diverse scripts. This evaluation utilizes the FLORES-101 benchmark~\citep{goyal2022flores}, and corresponding SentencePiece BLEU scores are summarized in Table~\ref{tab:flores}. Our findings show that {\textsc BLT} surpasses the baseline trained with the Llama 3 tokenizer, yielding an overall improvement of 2 BLEU points for translations into English and a 0.5-point improvement for translations from English. For high-resource language pairs, the translation quality of {\textsc BLT} matches or modestly exceeds that of Llama 3. Notably, {\textsc BLT} achieves stronger performance over Llama 3 across many pairs in lower-resource families, highlighting the utility of byte modeling approaches for improved generalization to rare or underrepresented byte sequences.

\subsection{Training {\textsc BLT} from Llama 3}
\label{sec:distillation}
We investigate a method by which {\textsc BLT} models can capitalize on pre-existing, pre-trained tokenizer-based architectures to improve training efficiency and performance. This approach involves initializing the global transformer parameters of {\textsc BLT} with those obtained from a pre-trained Llama 3.1 model. During the subsequent training phase, we update the global transformer's weights using a learning rate set to one-tenth that of the local encoder and decoder, maintaining the optimal number of training steps as specified for Llama 3. Our results, as shown in Table~\ref{tab:distill}, demonstrate that the {\textsc BLT} model initialized from Llama 3.1 notably surpasses both standard Llama 3 and vanilla {\textsc BLT} models, each trained for an equivalent computational budget (in terms of flops). Furthermore, when assessed against our {\textsc BLT}-Entropy variant—which used a considerably larger corpus (1T tokens), as documented in Table~\ref{tab:evals}—{\textsc BLT} from Llama 3.1 attains even higher scores on the MMLU benchmark. These findings indicate that this initialization strategy is a highly efficient means to decrease the overall training flops required while maintaining, or even improving, model quality.

This approach can be interpreted as adapting models that rely on tokenizers into ones that do not, thereby transforming a pre-trained LLaMA 3.1 model into a {\textsc BLT} model. For a thorough evaluation, Table \ref{tab:distill} presents the original LLaMA 3.1, which was pre-trained on 15T tokens, and compares its results with those of {\textsc BLT} models generated from LLaMA 3. While our model maintains similar performance on MMLU and HumanEval, we observe a more pronounced decrease on several other benchmarks. These findings indicate that additional research is required to better exploit the pre-trained model’s capabilities, with a focus on refining data composition and tuning other critical hyperparameters.

\section{Ablations and Discussion}
\label{section:ablations}
Here, we present ablation studies that support the design decisions behind the architecture of {\textsc BLT}, as well as the patching methodology and the selection of hyper-parameters for the 8B parameter version of {\textsc BLT} trained on the {\textsc BLT}-1T dataset.

\paragraph{Entropy Model Hyper-parameters} In order to analyze how entropy model capacity and the length of the context window influence scaling behavior, we experiment with byte-level entropy transformer models, adjusting the number of parameters from 1 million to 100 million and altering context window sizes from 64 up to 512. We generate scaling law plots of bits-per-byte (bpb) against training FLOPs for our $400m$ and $1b$ {\textsc BLT} models, each trained on the Llama-2 dataset; these plots are displayed in \autoref{fig:entropy_ablation}. Our observations indicate that scaling up either model size or context length improves performance, though the benefits taper off when model sizes exceed 50 million parameters.

\paragraph{Types of Patching}

We perform an ablation study on the four patching strategies outlined in Section~\ref{section:patching}: 1) Strided Patching, employing strides of 4 and 6, 2) Whitespace-based Patching, 3) Patch selection via BPE Tokenizer following the Llama 3 tokenizer, and 4) Entropy-driven patching utilizing a lightweight byte-level language model.

Although dynamic patching shortens the actual length of the sequences, we regulate sequence length so that every patching strategy operates under comparable context lengths. Across both training and inference phases, each model processes an equivalent expected number of bytes per sequence, thereby eliminating potential confounding effects linked to differences in context size. The findings from these comparative studies are presented in Figure~\ref{fig:scaling-compute-opt}. Every alternative patching approach surpasses static patching in performance, with space patching demonstrating results nearly equivalent to dynamic entropy based patching.

\autoref{tab:evals_ablation} displays comparative benchmark results for {\textsc BLT} models utilizing three different approaches: tokenizer-based models, space patching, and entropy-based patching, each trained for the ideal number of steps on the Llama 2 dataset~\citep{dubey2024llama}. While space patching constitutes a more straightforward approach—requiring no real-time execution of the entropy model during training—our results indicate that the improvements from entropy-based patching observed in the scaling analyses (see Section~\ref{section:scaling}) persist in downstream benchmark evaluations as well.\footnote{The space patching results shown stem from prior runs that excluded cross-attention; however, comparable trends are apparent when cross-attention is incorporated.}

\paragraph{Cross-Attention} 
In~\autoref{tab:crossattn_ablations_1b}, we systematically examine the impact of inserting cross-attention mechanisms at different stages within the encoder and decoder of {\textsc BLT}. For the encoder side, our experiments compare three initialization strategies for the cross-attention queries: (1) using an identical learned embedding for each global state; (2) employing a hash-based embedding derived from the patch's byte values; and (3) applying a pooling operation over the encoder layer's hidden states corresponding to the patch bytes.

Our results indicate that incorporating cross-attention within the \textit{decoder} yields the greatest benefit. Employing cross-attention in the encoder brings marginal gains, and these occur exclusively when queries are initialized through pooling. Furthermore, we observe that the advantage of cross-attention is more pronounced for the Common-Crawl dataset and becomes increasingly significant as patch sizes grow.

\paragraph{n-gram Hash Embeddings} 

We conduct ablation experiments using n-gram hash embedding vocabularies of 0, 100K, 200K, and 400K, and summarize the results in Table~\ref{tab:ngram_ablations_1b}. Hash embeddings consistently yield improvements across all domains, with the largest impact observed for Wikipedia and Github (an increase of 0.04 bpb, compared to a 0.01 bpb difference after 15k steps at 8B scale). At 8B scale, reducing the hash vocabulary from 500K to 300K affects performance by approximately 0.001 bpb after 15k steps. These results suggest that incorporating hashes is crucial for {\textsc BLT} to reach the performance levels of models using tokenizers; nevertheless, benefits plateau beyond 300K hashes. Furthermore, the performance gains from hashes and cross-attention mechanisms appear mostly additive, as each contributes improvements to distinct datasets.

\paragraph{Local Model Hyperparameters} Table \ref{tab:local_layers} presents an ablation study exploring different configurations for the number of layers utilized in the local encoder and decoder. Results indicate that when hash n-gram embeddings are used, {\textsc BLT} achieves strong performance even with a very minimal encoder—consisting of only a single layer—while a more substantial decoder proves beneficial.

\section{Related Work}
\label{section:related_work}

\textbf{Character-Level RNNs:} Since the advent of neural architectures, Character Language Modeling has garnered significant interest~\citep{sutskever2011generating,mikolov2012subword,graves2013generating}, primarily due to its innate capacity to represent out-of-vocabulary terms naturally, eliminating the need for traditional back-off strategies. For example, \cite{kim2016character} introduce a model that encodes sequences exclusively at the character level using convolutional and highway networks, which in turn interface with LSTM-driven RNNs. Their results demonstrate that the model achieves parity with leading RNN-based language models on English datasets and notably surpasses them on languages with rich morphology—a notable benefit of character-level LLMs. In a similar vein, \cite{kenter2018byte} employ byte-level LSTM networks for machine comprehension tasks and find that these models outperform their word-level counterparts, particularly on Turkish and Russian, both typified by complex morphology. Likewise, \cite{zhang2015character} propose character-level convolutional approaches for classification, observing that such models exceed the performance of word-level baselines in select scenarios. In pursuit of further enhancements, \citet{chung2019hierarchical} leverage hierarchical LSTM structures augmented with boundary detection mechanisms at various levels, allowing the model to autonomously uncover latent textual hierarchies and thereby boost character-level language modeling performance. Departing from attention-based methods, ByteNet~\cite{kalchbrenner2016neural} instead employs convolutional layers at the character level for neural machine translation tasks.

\textbf{Character-Level Transformers:} The introduction of attention-based transformer architectures~\citep{vaswani2017attention}, paired with subword tokenization techniques~\citep{sennrich-etal-2016-neural}, has led to substantial gains for neural approaches in language modeling and benchmark performance. Nonetheless, employing word or sub-word representations imposes an implicit inductive bias regarding the granularity at which such models function. In efforts to merge the advancements of transformers with the encouraging outcomes seen in early character-level language modeling, \citet{al2019character} experimented with training extremely deep transformer networks, augmented by auxiliary losses. This approach produced transformer models that surpassed earlier LSTM-based character language models. Despite these improvements, the resulting models still lagged noticeably behind leading word-level LLMs. Similarly, the creators of GPT-2~\citep{radfordgpt2} found that byte-level language models, when applied to extensive datasets like the 1 billion word benchmark, remained less effective compared to word-level alternatives.

\cite{Choe2019BridgingTG} showed that transformer-based byte-level LLMs can surpass subword-level LLMs of similar size, although these models demand substantially greater computational resources and require significantly more training time. Likewise, \cite{el2020characterbert} trained CharFormer, a BERT variant that derives word-level representations using convolutional layers over character embeddings, reporting performance gains in the medical domain, but with increased computational expenditure. In a related line of work, \cite{clark2022canine} introduced CANINE, an encoder-only model with 150M parameters that processes character sequences directly. CANINE's core comprises a deep transformer stack reminiscent of our global model, along with a local transformer module and strided convolutions to achieve input character downsampling; this design achieves superior results compared to its token-level counterpart, mBERT, on multilingual downstream benchmarks. ByT5~\citep{xue2022byt5} investigated methodologies for constructing encoder-decoder models at the byte level, deliberately omitting patching mechanisms. This model displayed enhanced resilience to input noise and demonstrated competitiveness with traditional tokenizer-based models despite being trained with one-fourth the data. However, the absence of patching necessitated costly attention computations over each byte, dramatically escalating computational demands. Shifting from subword to byte-level modeling inevitably prolongs input sequences, posing considerable efficiency challenges for scaling such models. In recent developments, leveraging the Mamba Architecture~\citep{gu2023mamba}, which preserves a fixed memory footprint across extensive context windows, \cite{wang2024mambabyte} constructed a byte-level Mamba-based model also operating without patching. Their results indicate outperformance over byte-level transformer models, assessed in a FLOP-controlled comparison at the 350M parameter scale based on bits-per-byte across several benchmark datasets.

\textbf{Patching-based approaches:} Leveraging patching methodologies can significantly reduce the typically large computational demands of byte-level LLMs, while potentially maintaining their effectiveness. Several studies have reported promising outcomes with patching methods at moderate scales in terms of both model parameters and training byte counts. For instance, \cite{nawrot-etal-2022-hierarchical} investigated static patching with both downsampling and upsampling operations, introducing the hourglass transformer, which achieved superior results compared to other byte-level baselines at the 150M parameter level. Building upon this, \cite{nawrot-etal-2023-efficient} enhanced patching efficiency through dynamic patching strategies, incorporating a learned boundary-predictor, a boundary-predictor trained with guidance from specific tokenizers, and an entropy-driven patching model akin to {\textsc BLT}. Their findings indicated that these dynamic schemes surpassed standard transformer models on language modeling tasks, specifically at the 40M parameter scale with about 400M training tokens. Meanwhile, \cite{lester2024training} explored training LLMs on sequences compressed via arithmetic coding to surpass the compression achievable by BPE. Employing an equal-info windows strategy, they managed to exceed the performance of byte-level baselines in language modeling, though these models still lagged behind those trained on subword units.

Our study is primarily motivated by and builds upon the foundational work of MegaByte~\citep{yu2023megabyte}, a decoder-only causal LLM architecture that implements static patching and concatenates feature representations to map bytes into patches. MegaByte subsequently employs a dedicated local model on the decoder to reconstruct bytes from these patches. Their results indicate that MegaByte achieves parity with tokenizer-based models at the 1B parameter level, evaluated on a dataset containing roughly 400B bytes. In our own experimental evaluation, we consistently ablate MegaByte, finding that its reliance on static patching causes it to underperform relative to state-of-the-art compute-optimal tokenizer models when measured in a flop-matched setting; we further demonstrate how {\textsc BLT} can close this performance gap. Similarly, \cite{slagle2024spacebyte} report comparable findings with respect to MegaByte and propose enhancing static patching by incorporating patches over whitespaces and space-like bytes as well as integrating a local encoder. Their results point to modest gains for byte-level approaches over tokenized transformer models under controlled compute for specific tasks, such as those found on Github and arXiv at a 1B parameter scale. In this work, we additionally evaluate this model variant and highlight that further architectural advancements are required if byte-level models are to be competitive with leading token-based architectures like Llama 3 at larger scales.

\section{Limitations and Future Work}

In this study, our approach to architectural decisions involves training models for the ideal number of steps based on the findings from Llama~3~\citep{dubey2024llama}. Nevertheless, it is important to note that these scaling laws were derived from experiments on BPE-level transformers, which may not yield the most effective ratios of data and parameter sizes for {\textsc BLT}. Developing scaling laws specifically tailored to {\textsc BLT} remains a direction for future research and may reveal more advantageous scaling behaviors for this architecture. Furthermore, the majority of our experiments were performed with models containing up to 1B parameters; it is plausible that the most effective architectural configurations could shift as the model size increases to 8B parameters and above, possibly resulting in enhanced performance at these larger scales.

Current transformer libraries and codebases are primarily optimized for architectures based on tokenizers. Although we provide theoretical flop-matched experimental results and leverage efficient tools like FlexAttention to support layers diverging from standard transformer designs, our implementations may still lag behind tokenizer-based models in terms of wall-clock performance and could be improved with additional optimizations.

Whereas {\textsc BLT} relies on a separately pre-trained entropy model for its patching approach, a promising avenue for future investigation is the joint end-to-end training of the patching mechanism itself. In Section \ref{sec:distillation}, we conduct preliminary experiments suggesting the potential success of converting tokenizer-based models like Llama 3—trained on over 10T tokens—into byte models by reusing their pretrained global transformer weights and keeping these weights fixed. Advancing this line of research could yield strategies that not only preserve the advantages brought by byte-ification, but also achieve superior performance relative to their tokenizer-based predecessors, all without necessitating training from the ground up.

\section{Conclusion}
\label{section:conclusion}

In this work, we introduce the Byte Latent Transformer (\textbf{BLT}), an innovative model architecture that challenges the standard reliance on fixed-vocabulary tokenization in large-scale language models. {\textsc BLT} employs a trainable and adaptive mechanism to cluster bytes into variable-length patches, thereby directing computational effort according to the intrinsic complexity of the data. This dynamic allocation enhances both computational efficiency and model robustness. Through comprehensive scaling experiments, we show that {\textsc BLT} attains performance on par with established tokenization-based architectures such as Llama 3 when trained at scales up to 8B parameters and 4T bytes, while offering up to 50\% fewer inference flops at only marginal costs to standard evaluation metrics. Additionally, {\textsc BLT} facilitates joint scaling of both the model size and patch size under a constant inference computational constraint, an approach well-suited to practical compute environments. By operating directly on raw byte sequences, {\textsc BLT} not only exhibits enhanced resilience to noisy or atypical data but also gains improved capabilities for modeling sub-word features. Taken together, these advances establish {\textsc BLT} as a compelling and adaptable alternative to classical tokenization-based models, capable of supporting more robust and efficient large language systems.

\end{document}